% source: https://github.com/pfradin/luadraw/discussions/245#discussioncomment-16870227 (pfradin, block 1)
\documentclass[border=2 mm]{standalone}
\usepackage[svgnames]{xcolor}
\usepackage[3d]{luadraw}
\usepackage{fouriernc}
\begin{document}
%%https://tex.stackexchange.com/questions/739936/fill-area-between-hyperbola-and-line-with-pgfplots
\begin{luadraw}{name=fill}
local ld = luadraw
local cpx = ld.cpx
local sqrt = math.sqrt
local Z, i = cpx.Z, cpx.I
local xmin, xmax, ymin, ymax = -1, 7, -3, 5
local a, b = 2,3
local g = ld.graph:new{window={xmin, xmax, ymin, ymax}}  
g:Daxes({0,1,1}, {grid=true, gridstyle = "dashed",arrows="-Stealth", legend={"$x$","$y$"}})
g:Linewidth(12)
local f = function(x) return 4*x-16  end
local h = function(x) return sqrt(x^2-9) end
local k = function(x) return -1*sqrt(x^2-9) end
g:Dinequalities({f,'>', h,'<'}, {view={3,5,0,5},draw_options="draw=none,fill=green,fill opacity=0.6"})
g:Dinequalities({f,'>', k,'>'}, {view={3,4,-2,0},draw_options="draw=none,fill=green,fill opacity=0.6"})
g:Dcartesian(f, {draw_options="violet"});
g:Dcartesian(h, {draw_options="blue, line cap=round"});
g:Dcartesian(k, {draw_options="blue, line cap=round"});  
g:Dlabel("$y=4x-16$",Z(6,1.5),{node_options="violet"})
g:Dlabel("$x^2-y^2=9$",Z(1.5,1.2),{node_options="violet"})
g:Show()
\end{luadraw}
\end{document}
